\section{Questions suggested by the comparison}\label{sec:questions}
The finite inequalities separate several issues that the leading constant alone cannot decide.

\paragraph{The $110$ and common-class constants.}
Do $\kappa_d(n)$ and $\kappa_c(n)$ converge? A natural question for $110$ avoidance is whether its limit is $C_+$, as in the $100$ class. A proof would need a $110$-preserving lower comparison with subexponential loss and a sufficiently small length shift. The doubled-seed construction spends an extra $r$ positions and currently gives $C_-$. It does not establish that $C_-$ is optimal for the common class, or that the two classes share a limit.

\paragraph{Sublinear terms and relative equivalents.}
What is the first correction after $C_+n$ for $a_n$ and $e_n$? The present upper error is $O(n(\log\log n)^2/\log n)$, while the lower comparison loses $O(n\log\log n/\log n)$. Determining their actual sublinear terms would clarify how much entropy the ordinary words gain or lose relative to triangular arrays. Even equality through several logarithmic orders would not imply $a_n\sim b_n$ or $e_n\sim b_n$: a relative equivalent requires logarithmic error tending to zero. The ratio $e_n/a_n$, which measures the cost of excluding a third occurrence, is also left open.

\paragraph{Sharper skeleton and fiber counts.}
The upper bound treats marked positions, marked labels, and boundaries independently, although many resulting skeletons are inconsistent. Can their compatibility be counted sharply? On the lower side, $2^r$, $\binom{p-1}{r-1}$, and $Q+1$ are worst-case fiber bounds. Their typical values may be much smaller. A refined analysis could reduce the error without pretending the undecorated maps are injective.

\paragraph{Typical word structure.}
The triangular lower construction uses $r\sim2n/\log n$, but its seed size is a construction parameter, not automatically the number of repeated occurrences in a uniform ordinary avoider. Likewise, an upper-bound saddle does not prove concentration of that statistic. It remains to determine typical multiplicities, the number and lengths of maximal increasing runs, and the prevalence of merged block boundaries for these word classes. Such results require distributional comparisons stronger than total-count inequalities.

\paragraph{More accurate threshold inversion.}
The present inverse has a diverging absolute error. A tighter count remainder would transport to a tighter inverse by division by a slope of order $\log n$, but a uniformly exact rounding rule needs far more: a sufficiently accurate count approximation and control near integer boundaries. For $d,c$, determining the linear constant is already necessary to identify the second inverse coefficient. No finite numerical fit or rounding of \eqref{eq:inverseexact} resolves these questions.
